\documentclass[../DoAn.tex]{subfiles}
\begin{document}

\begin{center}
    \Large{\textbf{ABSTRACT}}\\
\end{center}
\vspace{1cm}
\hspace*{1.5cm}In the current educational landscape, training management and online exam monitoring face significant limitations. Traditional attendance methods are time-consuming and prone to proxy attendance. Furthermore, popular Learning Management Systems (LMS) like Moodle lack reliable biometric verification support, while commercial proctoring software is expensive and requires substantial server-side GPU resources to operate.

\hspace*{1.5cm}To address these challenges, this graduation thesis proposes a comprehensive and synchronized software ecosystem entitled "Building software to optimize the attendance and student management process". The system is developed with a React-based Web dashboard and a React Native mobile application. The entire platform operates on a microservice architecture powered by a Spring Boot backend, a PostgreSQL database, and centralized authentication secured by Keycloak SSO.

\hspace*{1.5cm}The key technical contributions of this thesis include three core modules. First, a group attendance solution (Group FaceID) utilizing RetinaFace and ArcFace models to process classroom photos and identify students within 1.5 to 2.5 seconds. Second, a two-layer geofencing mechanism that detects and rejects check-in requests using simulated GPS coordinates by validating OS-level Mock Location flags on mobile devices. Finally, a real-time online exam monitoring (AI Proctoring) engine leveraging MediaPipe FaceMesh, 3D solvePnP pose estimation, and Gaze Tracking streaming over WebSockets with a latency under 80 ms on standard CPUs.

\hspace*{1.5cm}Experimental evaluations involving 150 students and 10 lecturers demonstrate outstanding stability and performance. The system facilitates group attendance verification, logs and blocks simulated GPS coordinates when Mock Location is active, and maintains API response times below 150 ms. These results confirm the high practicality and feasibility of the proposed solution in optimizing modern educational management processes.

\end{document}